% !TEX root = ../paper_P5_minreport.tex
% Section — Iter 201 P5 task-slice stratified algorithm-vs-stack variance ratio
\subsection{Task-stratified ratio: iter-193's headline is task-conditional (iter 201)}
\label{sec:p5-iter201-task-stratified}

Section~\ref{sec:p5-iter193-ratio} reported a single corpus-wide
stack-to-label variance ratio of $60.6\times$ on reward,
$10.3\times$ on zvf, and $4.8\times$ on completion length, computed
on the $98$ \texttt{mega\_20260704} cells treated as one bag.
Companion work (iter-197) robustness-audited that headline at the
corpus level with paired bootstrap, worst-axis stress, jackknife,
and a $5$-axis composite. Both treated the corpus as a single
homogeneous pool and reported a single ratio per channel.

We close the deeper question here: does the iter-193 headline
\emph{survive stratification by task slice}? Stratifying by
\texttt{gsm8k\_easy}, \texttt{gsm8k\_hard}, and
\texttt{humaneval\_subset} and recomputing the ratio within each
task slice reveals a sharp task-conditional qualification.

\begin{table}[t]
\centering\small
\caption{Per-(task slice, channel) stack-vs-label variance ratio with
$95\%$ stratified bootstrap CI. The iter-193 corpus-wide headline
($60.6\times$ reward) is the GSM8K-only signal; on
\texttt{humaneval\_subset} the reward and zvf ratios are
\emph{structurally zero} because every cell in that slice has
\texttt{mean\_reward=0.0} and \texttt{zvf=1.0} (all-zero rewards
$\Rightarrow$ SS$_{\text{total}}=0$). Only the completion-length
channel has variance on every task slice and produces a meaningful
ratio on all three ($34$--$73\times$).}
\label{tab:p5-iter201-per-task}
\begin{tabular}{llccrc}
\toprule
task slice & channel & top axis & top $\eta^2$ & ratio & 95\% CI \\
\midrule
\texttt{gsm8k\_easy}      & zvf    & $G$             & $0.887$ & $19.5\times$  & $[6.0, 85.2]$ \\
\texttt{gsm8k\_hard}      & zvf    & $G$             & $0.641$ & $14.1\times$  & $[4.5, 66.8]$ \\
\texttt{humaneval\_sub}   & zvf    & model\_family   & $0.000$ & $0.0\times$   & $[0.0, 0.0]$ \\
\midrule
\texttt{gsm8k\_easy}      & reward & model\_family   & $0.993$ & $133.0\times$ & $[12.6, 548.5]$ \\
\texttt{gsm8k\_hard}      & reward & model\_family   & $0.991$ & $132.8\times$ & $[12.8, 573.7]$ \\
\texttt{humaneval\_sub}   & reward & model\_family   & $0.000$ & $0.0\times$   & $[0.0, 0.0]$ \\
\midrule
\texttt{gsm8k\_easy}      & len    & model\_family   & $0.382$ & $51.1\times$  & $[3.4, 228.5]$ \\
\texttt{gsm8k\_hard}      & len    & temperature     & $0.256$ & $34.3\times$  & $[2.2, 152.6]$ \\
\texttt{humaneval\_sub}   & len    & model\_family   & $0.547$ & $73.2\times$  & $[6.7, 312.5]$ \\
\bottomrule
\end{tabular}
\end{table}

\textbf{Findings.}
\textit{(i)} \textbf{Headline task-conditional qualification.} The
iter-193 corpus-wide $60.6\times$ reward ratio is the GSM8K average;
on \texttt{humaneval\_subset} the ratio is structurally zero because
all $34$ cells have \texttt{mean\_reward}$=0.0$ and
\texttt{zvf}$=1.0$ (zero reward variance $\Rightarrow$
$\eta^2_{\text{stack}}=\eta^2_{\text{algo}}=0$). MIN-REPORT's
``Report the Stack'' thesis is trivially true on degenerate tasks
(no signal $\Rightarrow$ no variance $\Rightarrow$ no axis can
dominate). The $60.6\times$ headline applies \emph{only} to tasks
where models achieve non-zero success rate. \textit{(ii)}
\textbf{Completion length is task-robust.} Only the len channel has
variance on every task slice and produces a meaningful ratio on all
three ($34$--$73\times$). When reward is degenerate (early training
or saturated baseline), len is the recommended fallback for
MIN-REPORT stack validation. \textit{(iii)} \textbf{$G$ dominates
zvf within GSM8K} ($\eta^2=0.887$ on \texttt{gsm8k\_easy}, $0.641$
on \texttt{gsm8k\_hard}) --- the iter-193 finding is robust
\emph{within GSM8K} but not corpus-wide. \textit{(iv)}
\textbf{Model family dominates reward everywhere it is defined}
($\eta^2\ge 0.99$ on both GSM8K slices); this is the strongest
single-task pattern in the corpus. \textit{(v)} \textbf{Five
hypotheses}: H1 (point ratio $>1$ on every cell) FAIL on
\texttt{humaneval\_subset} --- the sharp qualification. H2 (CI
excludes 1 on $\ge 2/3$ tasks per channel) PASS. H3 (between-task
ratio variance $>0.5$) PASS --- variance is $101$--$5885$ per
channel. H4 (top axis consistent on $\ge 2/3$ tasks per channel)
PASS. H5 (zvf top axis varies across tasks) PASS.

\textbf{Implication for MIN-REPORT.} The iter-193 headline is
preserved for the GSM8K subset (where models actually train) but
qualified corpus-wide: on degenerate tasks the stack-vs-label
distinction collapses to a 0/0 form. Practitioners should report
per-task ratios, not corpus-wide averages, and should validate the
stack-vs-label dominance on \emph{at least one task where the
algorithm achieves non-zero success rate} before claiming the stack
is the dominant driver.